%! Rational Functions and Vertical Asymptotes — Unit 3, Lesson 3.3 — Slides (Beamer)
\documentclass[aspectratio=169, 11pt]{beamer}
\usepackage{precalculus-beamer}   % provides \navyheader{} and \sectionlabel[color]{}

\begin{document}

% TITLE SLIDE
{
\setbeamercolor{background canvas}{bg=navy}
\begin{frame}[plain]
  \vspace{2.8cm}\hspace{0.55cm}%
  \begin{minipage}{0.9\textwidth}
    {\color{goldacc}\bfseries\small LESSON 1.9}\\[0.5cm]
    {\color{white}\bfseries\LARGE Rational Functions and Vertical Asymptotes}\\[1.2cm]
    {\color{skymid}\small Precalculus~$\cdot$~Unit 3: Rational Functions}
  \end{minipage}
\end{frame}
}

% HOOK SLIDE
\begin{frame}[t, plain]
  \navyheader{Hook: Bacterial Concentration Near a Drain}
  \sectionlabel[goldacc]{Think About It}

  Concentration is modeled by $C(x) = \dfrac{1}{x - 3}$, where $x$ is distance (cm) from the drain.

  \bigskip
  \begin{columns}[T]
    \column{0.52\textwidth}
      \begin{block}{Table of Values Near $x = 3$}
        \renewcommand{\arraystretch}{1.3}
        \begin{tabular}{|c|c|}
        \hline
        $x$    & $C(x)$ \\
        \hline
        2.9    & $-10$ \\
        2.99   & $-100$ \\
        2.999  & $-1{,}000$ \\
        3.001  & $1{,}000$ \\
        3.01   & $100$ \\
        3.1    & $10$ \\
        \hline
        \end{tabular}
      \end{block}

    \column{0.44\textwidth}
      \begin{block}{Discussion Questions}
        \begin{itemize}\setlength{\itemsep}{6pt}
          \item Can we evaluate $C(3)$? Why or why not?
          \item What happens to $C(x)$ as $x$ gets closer to 3?
          \item Why does the sign of $C(x)$ change as we cross $x = 3$?
        \end{itemize}
      \end{block}
  \end{columns}
\end{frame}

% VERTICAL ASYMPTOTES: RULE + EXAMPLE
\begin{frame}[t, plain]
  \navyheader{Vertical Asymptotes}

  \begin{block}{Rule (EK 1.9.A.1)}
    $r(x) = p(x)/q(x)$ has a \textbf{vertical asymptote at $x = a$} when:
    \begin{itemize}\setlength{\itemsep}{4pt}
      \item $q(a) = 0$ \quad (denominator is zero at $a$), \textbf{AND}
      \item $p(a) \ne 0$ \quad (numerator is \emph{not} zero at $a$).
    \end{itemize}
    \textit{Multiplicity rule:} VA also occurs if the multiplicity of $a$ in the denominator
    exceeds its multiplicity in the numerator.
  \end{block}

  \bigskip
  \begin{columns}[T]
    \column{0.48\textwidth}
      \textbf{Example 1:} $r(x) = \dfrac{x+2}{x-3}$
      \begin{itemize}\setlength{\itemsep}{5pt}
        \item Denominator zero: $x = 3$
        \item Numerator at $x=3$: $5 \ne 0$
        \item \textbf{VA at $x = 3$}
      \end{itemize}

    \column{0.48\textwidth}
      \textbf{Example 2:} $r(x) = \dfrac{x-1}{(x-1)^2(x+2)}$
      \begin{itemize}\setlength{\itemsep}{5pt}
        \item $x=1$: denom.\ mult.\ 2 $>$ numer.\ mult.\ 1 $\Rightarrow$ \textbf{VA}
        \item $x=-2$: numer.\ $= -3 \ne 0$ $\Rightarrow$ \textbf{VA}
      \end{itemize}
  \end{columns}
\end{frame}

% ONE-SIDED LIMITS
\begin{frame}[t, plain]
  \navyheader{One-Sided Limits Near a Vertical Asymptote}

  For $r(x) = \dfrac{1}{x - 2}$:

  \bigskip
  \begin{columns}[T]
    \column{0.50\textwidth}
      \begin{block}{Table of Values}
        \renewcommand{\arraystretch}{1.3}
        \begin{tabular}{|c|c|}
        \hline
        $x$    & $r(x)$ \\
        \hline
        1.9    & $-10$ \\
        1.99   & $-100$ \\
        1.999  & $-1{,}000$ \\
        2.001  & $1{,}000$ \\
        2.01   & $100$ \\
        2.1    & $10$ \\
        \hline
        \end{tabular}
      \end{block}

    \column{0.46\textwidth}
      \begin{block}{Limit Notation (EK 1.9.A.2)}
        \[
          \lim_{x \to 2^-} r(x) = -\infty
        \]
        \[
          \lim_{x \to 2^+} r(x) = +\infty
        \]
        \medskip
        \textit{Sign analysis:} when $x > 2$, $(x-2) > 0$, so $r(x) > 0 \to +\infty$.\\
        When $x < 2$, $(x-2) < 0$, so $r(x) < 0 \to -\infty$.
      \end{block}
  \end{columns}
\end{frame}

% CLASS PRACTICE
\begin{frame}[t, plain]
  \navyheader{Class Practice}
  \sectionlabel[redacc]{Try It --- With a Partner}

  \begin{enumerate}\setlength{\itemsep}{10pt}
    \item For $r(x) = \dfrac{x + 1}{x^2 - 4}$:
          \begin{itemize}\setlength{\itemsep}{4pt}
            \item Factor the denominator.
            \item Find the denominator zeros.
            \item Check the numerator at each zero.
            \item State all vertical asymptotes.
          \end{itemize}

    \item For $r(x) = \dfrac{3}{(x - 5)^2}$:
          \begin{itemize}\setlength{\itemsep}{4pt}
            \item Identify the vertical asymptote.
            \item Use sign analysis: what is the sign of $(x-5)^2$ near $x=5$?
            \item Write $\displaystyle\lim_{x \to 5^+} r(x)$ and $\displaystyle\lim_{x \to 5^-} r(x)$.
          \end{itemize}
  \end{enumerate}
\end{frame}

\end{document}
